\begin{trapbox}
\begin{description}[style=nextline, leftmargin=2.8em, labelsep=0.5em, itemsep=0.35em]
    \item[\textbf{\textcolor{CTrap}{易错点一（误判条件）}}] 认为任意四个下标都成立，忽略 $m+n=p+q$。
    \item[\textbf{\textcolor{CTrap}{正确理解（条件必须满足）：}}] 只有当下标和相等时才能直接应用结论，否则不能直接用。
    \item[\textbf{\textcolor{CTrap}{错因分析（忽视等量关系）：}}] 将结论错误记忆为“任意两项乘积相等”，未检查下标和。
\end{description}

\medskip
\begin{description}[style=nextline, leftmargin=2.8em, labelsep=0.5em, itemsep=0.35em]
    \item[\textbf{\textcolor{CTrap}{易错点二（正负漏解）}}] 由乘积求单项时只取正值，忽略负数情况。
    \item[\textbf{\textcolor{CTrap}{正确理解（考虑公比符号）：}}] 等比数列公比可正可负，开平方后应得 $\pm$。
    \item[\textbf{\textcolor{CTrap}{错因分析（惯性思维）：}}] 受等差数列影响，默认所有项为正；实际等比数列可正负交替。
\end{description}

\medskip
\begin{description}[style=nextline, leftmargin=2.8em, labelsep=0.5em, itemsep=0.35em]
    \item[\textbf{\textcolor{CTrap}{易错点三（混淆等差等比）}}] 将等比性质与等差性质混淆，认为是 $a_m+a_n=a_p+a_q$。
    \item[\textbf{\textcolor{CTrap}{正确理解（乘法与加法区别）：}}] 等比数列是乘积相等，等差数列是和相等。
    \item[\textbf{\textcolor{CTrap}{错因分析（记忆混淆）：}}] 未区分两种数列的运算本质，记反规律。
\end{description}
\end{trapbox}
